\input{preamble.tex}

\title{Отчет по лабораторной работе № 12. \\ Настройка NAT (Network Address Translation)}
\author{Сущенко Алина Николаевна\\ НПИбд-01-23}
\date{2026}

\begin{document}

\maketitle
\newpage

\tableofcontents

\newpage
\section{Цель работы}
Настроить NAT (Network Address Translation) для обеспечения доступа локальных сетей организации во внешнюю сеть (Интернет) и организовать доступ из внешней сети к внутренним серверам (WWW, почтовый, файловый, RDP).

\newpage
\section{Выполнение работы}

\subsection{Первоначальная настройка маршрутизатора ansusenko-gw-1}
Выполнили первоначальную настройку маршрутизатора провайдера: задали имя устройства, настроили пароли для доступа по VTY и консоли, установили пароль enable secret, включили шифрование паролей и создали пользователя admin (Рис.~\ref{fig:01}):

\begin{figure}[htbp]
    \centering
    \includegraphics[width=\textwidth]{../images/01.png}
    \caption{Первоначальная настройка маршрутизатора ansusenko-gw-1.}
    \label{fig:01}
\end{figure}

\subsection{Первоначальная настройка коммутатора ansusenko-sw-1}
Выполнили первоначальную настройку коммутатора провайдера схожим способом с маршрутизатором (Рис.~\ref{fig:02}):

\begin{figure}[htbp]
    \centering
    \includegraphics[width=\textwidth]{../images/02.png}
    \caption{Первоначальная настройка коммутатора ansusenko-sw-1.}
    \label{fig:02}
\end{figure}

\subsection{Настройка интерфейсов маршрутизатора ansusenko-gw-1}
Настроили интерфейсы маршрутизатора провайдера:
\begin{itemize}
    \item Интерфейс f0/0 - включили (no shutdown);
    \item Создали сабинтерфейс f0/0.4 с инкапсуляцией dot1Q для VLAN 4, назначили IP-адрес 198.51.100.1/28;
    \item Интерфейс f0/1 - включили, назначили IP-адрес 192.0.2.1/24 (возникла ошибка перекрытия с f0/0, но продолжили настройку) (Рис.~\ref{fig:03}):
\end{itemize}

\begin{figure}[htbp]
    \centering
    \includegraphics[width=\textwidth]{../images/03.png}
    \caption{Настройка интерфейсов маршрутизатора ansusenko-gw-1.}
    \label{fig:03}
\end{figure}

\subsection{Настройка интерфейсов коммутатора ansusenko-sw-1}
Настроили интерфейсы коммутатора провайдера:
\begin{itemize}
    \item Интерфейсы f0/1 и f0/2 настроили как транковые (switchport mode trunk);
    \item Создали VLAN 4 с именем \texttt{nat};
    \item Настроили интерфейс VLAN 4 (SVI) и включили его (no shutdown) (Рис.~\ref{fig:04}):
\end{itemize}

\begin{figure}[htbp]
    \centering
    \includegraphics[width=\textwidth]{../images/04.png}
    \caption{Настройка интерфейсов коммутатора ansusenko-sw-1.}
    \label{fig:04}
\end{figure}

\subsection{Настройка интерфейсов маршрутизатора msk-ansusenko-gw-1}
Настроили маршрутизатор сети «Донская»:
\begin{itemize}
    \item Интерфейс f0/1 - включили (no shutdown);
    \item Создали сабинтерфейс f0/1.4 с инкапсуляцией dot1Q для VLAN 4, назначили IP-адрес 198.51.100.2/28;
    \item Настроили маршрут по умолчанию (default gateway) через 198.51.100.1 (Рис.~\ref{fig:05}):
\end{itemize}

\begin{figure}[htbp]
    \centering
    \includegraphics[width=\textwidth]{../images/05.png}
    \caption{Настройка интерфейсов маршрутизатора msk-ansusenko-gw-1 и маршрута по умолчанию.}
    \label{fig:05}
\end{figure}

\subsection{Настройка пула адресов для NAT}
Настроили пул внешних IP-адресов \texttt{main-pool} с диапазоном 198.51.100.2 - 198.51.100.14 и маской 255.255.255.240 (Рис.~\ref{fig:06}):

\begin{figure}[htbp]
    \centering
    \includegraphics[width=\textwidth]{../images/06.png}
    \caption{Настройка пула адресов для NAT.}
    \label{fig:06}
\end{figure}

\subsection{Настройка списка доступа для NAT}
Теперь настроим список доступа для NAT (Рис.~\ref{fig:07}):

\begin{figure}[htbp]
    \centering
    \includegraphics[width=\textwidth]{../images/07.png}
    \caption{Настройка списка доступа для NAT.}
    \label{fig:07}
\end{figure}

\subsection{Настройка списка доступа для NAT (продолжение)}
Создали расширенный список доступа \texttt{nat-inet} с правилами доступа в Интернет (Рис.~\ref{fig:08}, \ref{fig:09}, \ref{fig:10}, \ref{fig:11}):

\subsubsection*{Доступ для сети «Донская»}
Разрешили доступ сети 10.128.3.0/24 к web-серверам yandex.ru (192.0.2.11) и stud.rudn.university (192.0.2.12) по порту TCP 80 (Рис.~\ref{fig:08}):

\begin{figure}[htbp]
    \centering
    \includegraphics[width=\textwidth]{../images/08.png}
    \caption{Правила доступа для сети «Донская» (dk).}
    \label{fig:08}
\end{figure}

\subsubsection*{Доступ для сети кафедры}
Разрешили доступ сети 10.128.4.0/24 к web-серверу esystem.pfur.ru (192.0.2.13) по порту TCP 80 (Рис.~\ref{fig:09}):

\begin{figure}[htbp]
    \centering
    \includegraphics[width=\textwidth]{../images/09.png}
    \caption{Правила доступа для сети кафедры.}
    \label{fig:09}
\end{figure}

\subsubsection*{Доступ для сети администрации}
Разрешили доступ сети 10.128.5.0/24 к web-серверу rudn.ru (192.0.2.14) по порту TCP 80 (Рис.~\ref{fig:10}):

\begin{figure}[htbp]
    \centering
    \includegraphics[width=\textwidth]{../images/10.png}
    \caption{Правила доступа для сети администрации.}
    \label{fig:10}
\end{figure}

\subsubsection*{Доступ для компьютера администратора}
Разрешили полный доступ в Интернет компьютеру администратора с IP-адресом 10.128.6.200 (Рис.~\ref{fig:11}):

\begin{figure}[htbp]
    \centering
    \includegraphics[width=\textwidth]{../images/11.png}
    \caption{Правила доступа для компьютера администратора.}
    \label{fig:11}
\end{figure}

\subsection{Настройка NAT}
\subsubsection*{Port Address Translation (PAT) и настройка интерфейсов для NAT}
Настроили перегрузку PAT, связав список доступа \texttt{nat-inet} с пулом адресов \texttt{main-pool} с ключевым словом \texttt{overload} (Рис.~\ref{fig:12}):

\begin{figure}[htbp]
    \centering
    \includegraphics[width=\textwidth]{../images/12.png}
    \caption{Настройка PAT и интерфейсов для NAT.}
    \label{fig:12}
\end{figure}

\subsection{Настройка доступа из Интернета}

\subsubsection*{WWW-сервер}
Настроили статическую трансляцию порта TCP 80 для сервера 10.128.0.2 на внешний адрес 198.51.100.2 порт 80 (Рис.~\ref{fig:13}):

\begin{figure}[htbp]
    \centering
    \includegraphics[width=\textwidth]{../images/13.png}
    \caption{Статический NAT для WWW-сервера.}
    \label{fig:13}
\end{figure}

\subsubsection*{Файловый сервер}
Настроили статическую трансляцию порта TCP 20 и 21 (Рис.~\ref{fig:14}):

\begin{figure}[htbp]
    \centering
    \includegraphics[width=\textwidth]{../images/14.png}
    \caption{Статический NAT для файлового сервера.}
    \label{fig:14}
\end{figure}

\subsubsection*{Почтовый сервер}
Настроили статическую трансляцию портов TCP 25 (SMTP) и 110 (POP3) для сервера 10.128.0.4 на внешние адреса 198.51.100.4 порты 25 и 110 соответственно (Рис.~\ref{fig:15}):

\begin{figure}[htbp]
    \centering
    \includegraphics[width=\textwidth]{../images/15.png}
    \caption{Статический NAT для почтового сервера.}
    \label{fig:15}
\end{figure}

\subsubsection*{Доступ по RDP для компьютера администратора}
Настроили статическую трансляцию порта TCP 3389 (RDP) для компьютера администратора 10.128.6.200 на внешний адрес 198.51.100.10 порт 3389 (Рис.~\ref{fig:16}):

\begin{figure}[htbp]
    \centering
    \includegraphics[width=\textwidth]{../images/16.png}
    \caption{Статический NAT для RDP-доступа к компьютеру администратора.}
    \label{fig:16}
\end{figure}

\newpage
\section{Вывод}
В ходе лабораторной работы нам удалось настроить NAT (Network Address Translation) для обеспечения доступа локальных сетей организации во внешнюю сеть (Интернет) и организовать доступ из внешней сети к внутренним серверам (WWW, почтовый, файловый, RDP).

\end{document}
